\chapter{CAPITOLO 9}\label{capitolo-9}

\hfill\break

Mi misero in prigione nell'unica base kristang del pianeta, gestita
dalla trentina di kristang che erano stati vaccinati con un farmaco
sperimentale per proteggersi dal rischio biologico su Paradiso, tutti
volontari. Un'avanguardia dei più accaniti, o di quelli più
disperatamente alla ricerca di una promozione. Ognuno di loro era uno
stronzo fanatico. Ero tenuto in isolamento, in vero isolamento, non
vedevo nemmeno le lucertole mentre ero nella mia cella, il cibo veniva
consegnato attraverso un cassetto scorrevole una volta al giorno. Per
passare il tempo, guardavo la luce del sole cambiare direzione dalla
piccola finestra, in alto sul muro. E ascoltavo le urla che si
ripetevano a intermittenza e i colpi di fucile al mattino. Il quinto,
forse sesto giorno, ricevetti una visita dal corpo dei Marine degli
Stati Uniti.

«Bishop, sono il maggiore Cochrane, come stai, figliolo?» Si guardò
attorno nella cella, alla ricerca di un posto dove appendere il
cappello, poi se l'infilò sotto il braccio. Sembrava stesse anche peggio
di me, scarno e con occhiaie e borse sotto gli occhi. Gli ufficiali del
quartier generale dell'Unef stavano dando il buon esempio, tagliandosi
le razioni di cibo. Facevamo a gara per evitare di finire le scorte
prima che i nostri raccolti inaugurali fossero pronti.

Mi appoggiai contro la parete, non essendoci posto per sedersi, o
stendersi, a parte il duro, freddo pavimento. «Abbastanza bene, signore,
nei limiti del possibile. Niente di cui lamentarsi. Il cibo potrebbe
essere migliore e questo letto è duro, ma almeno non è grumoso»,
scherzai con fatica.

«Fanno sul serio», si accigliò Cochrane, «i kristang vogliono fare di te
un esempio, di te e degli altri che hanno sfidato i loro ordini.»

«Quanti altri? Ho sentito le squadre di fuoco, al mattino.»

Cochrane sembrava scosso. «Americani, che io sappia, diciassette, stiamo
cercando di ottenere un computo esatto. Più inglesi, indiani e cinesi e
un paio di soldati francesi.»

«Tutti nel braccio della morte?»

Annuì lentamente. «Il che non tiene in conto i morti di quando i
kristang si sono spazientiti e hanno colpito le unità dall'orbita. Danno
collaterale, dicono», e lanciò uno sguardo nervoso al soffitto, come se
fosse l'unico posto in cui le lucertole avrebbero potuto piazzare delle
cimici. «Si dà il caso che i siti che hanno colpito siano solo quelli in
cui i nostri esitavano a eseguire gli ordini di rappresaglia.» Mi guardò
con la coda dell'occhio. «I tuoi sono stati fortunati. Se non fossi
stato celebre tra i kristang, avrebbero spazzato via la tua squadra con
un attacco di cannone a rotaia, insieme a Habitrail. A quanto pare,
volevano che vedessi distruggere quella scuola. I kristang non tollerano
che specie inferiori li sfidino.»

«Per rappresaglia si intende quando si colpiscono soggetti che ti hanno
colpito. Le donne e i bambini criceti non ci hanno colpiti. Questo è
puro e semplice omicidio. Mi hanno fatto colonnello», indicai l'aquila
sulla mia uniforme, «hanno detto che ero un eroe per avere ucciso dei
soldati ruhar, ora vogliono uccidere me per aver rifiutato di
assassinare donne e bambini ruhar.»

«Hai fatto bene un sacco di cose, ma\ldots»

«Sì, hai fottuto una pecora.»

«Cosa?», chiese Cochrane, con gli occhi spalancati.

«È una battuta. La dimentichi. Potrei dire: ``Una sola lavata di capo
spazza via un centinaio di encomi'', eh?»

«Qualcosa di simile», annuì lui.

«Quindi, lei è qui per cosa, ascoltare la mia confessione, farmi da
avvocato in divisa, darmi un cucchiaio, così posso scavare un tunnel per
uscire di qui?» Indicai il pavimento di cemento indurito, o qualsiasi
sostanza dall'incredibile resistenza i kristang usassero come materiale
da costruzione.

Scosse la testa. «Mi dispiace dire che sono qui soltanto per dimostrarti
il supporto e l'interessamento dell'esercito degli Stati Uniti. I
kristang non hanno bisogno di una confessione e tu non hai bisogno di un
avvocato, perché non ci sarà nessun genere di processo o audizione.»

«Giusto, perché darsi la pena di fare giustizia quando si possono
giustiziare le persone?» Agitai la mano, provando dispiacere per
quell'uomo, che doveva trovarsi in una posizione incredibilmente
imbarazzante. «Risparmi il fiato. Sta per dire qualcosa su come gli
alieni abbiano diversi concetti di giustizia e disciplina militare. So
tutto questo. So anche che quest'idea di cosa sia giusto e cosa sia
sbagliato non dipende dagli occhi di chi guarda. È sbagliata essa
stessa.»

Cochrane alzò gli occhi alla minuscola finestra. Non c'era nulla che
potesse dire.

«Maggiore, non ha la sensazione che siamo dalla parte sbagliata di
questa guerra? Più imparo sulle lucertole, più sono convinto che stiamo
combattendo dalla parte nazista del conflitto. Durante la seconda guerra
mondiale, le Ss allinearono i civili e li fucilarono per rappresaglia,
per via dei loro attacchi ai soldati tedeschi. Gli Alleati non lo
fecero.» Indicai la bandiera degli Stati Uniti sulla mia uniforme: «Noi
siamo civilizzati. I nazisti non lo erano».

«I kristang non hanno attaccato la Terra», puntualizzò il mio
interlocutore.

«Maggiore, a meno che tu non avessi due belle fette di prosciutto su
quei cazzo di occhi\ldots»

Cochrane s'irrigidì, per rabbia, ma potevo vedere molta paura: «Attento
a come parli, soldato».

«Sono ancora di grado superiore a te, maggiore», enfatizzai l'ultima
parola. «L'esercito ha per caso revocato il mio rango teatrale
abbassandomi a quello di sergente? No? Allora sono ancora un colonnello
a tutti gli effetti e, a meno che non avessi due belle fette di
prosciutto sugli occhi», omisi la volgarità, perché non volevo che se ne
andasse, essendo l'unico umano che vedevo da giorni, «i ruhar hanno
razziato la Terra, ci hanno colpito per deteriorare la nostra capacità
industriale, così non saremmo più stati così utili alle lucertole. Non
hanno colpito le città, non hanno nemmeno colpito le basi militari,
hanno colpito soltanto infrastrutture industriali, centrali elettriche e
raffinerie. Se i ruhar non ci avessero colpito per primi, le lucertole
sarebbero apparse in cielo, ci avrebbero detto che ora lavoravamo per
loro e poi avrebbero spazzato via un paio di centinaia di migliaia di
persone per esprimere il loro punto di vista. L'unico motivo per cui non
ci siamo resi conto fin dall'inizio di essere ormai schiavi dei
kristang, è che eravamo loro tanto grati per aver cacciato via i ruhar.
I criceti hanno fatto un favore ai kristang. Noi umani non contiamo un
cazzo per entrambe le parti. Hai sentito le voci su quello che sta
succedendo a casa? Dai biscotti della fortuna? Shanghai? Parigi? San
Francisco?»

«Voci», protestò Cochrane a fatica.

«Voci a cui credo, molto più di quanto non creda alle stronzate
censurate che i nostri presunti alleati permettono al quartier generale
dell'Unef di rifilarci. Cosa mi aspetta? Sono qui da giorni.»

Cochrane distolse lo sguardo, poi incontrò il mio: «Plotone di
esecuzione, domani mattina, per i restanti prigionieri maschi».

«Prigionieri maschi?»

Emise una lunga espirazione. «I prigionieri maschi vengono semplicemente
fucilati. Le donne, loro\ldots{} conosci l'atteggiamento dei kristang
nei confronti delle femmine. Una donna che abbia autorità superiore a
quella degli uomini è già un male, ma una donna che sfidi gli ordini
impartiti dagli uomini? È una cosa che i kristang non sopportano, li
manda fuori di testa, e stanno impartendo una punizione esemplare. Le
donne prigioniere vengono spogliate, torturate e impiccate. Lentamente.»
Sembrava che fosse sul punto di vomitare. «Ci hanno costretti a
guardare, hanno costretto il comando dell'Unef a guardare, ieri.»

«Oh, merda. E non avete fatto niente?»

«Fare cosa? Colonnello», disse con enfasi, «di certo ammetterai che il
comando dell'Unef ha poca scelta. Siamo al termine della linea di
rifornimento più lunga della storia. Tutte le nostre munizioni,
forniture mediche, tutto il nostro cibo, deve arrivare qui su
un'astronave kristang. Al massimo, avevamo scorte di cibo sufficienti
per quattordici settimane, che erano tante, quando si prevedeva che
l'evacuazione dei ruhar sarebbe stata ultimata nel giro di un anno e le
navi di rifornimento dei kristang arrivavano qui puntuali come un
orologio svizzero. Quando gli assalti dei ruhar hanno interrotto le
consegne, ci è rimasto cibo per un mese circa. L'intera Expeditionary
Force potrebbe essere distrutta solo trattenendo scorte di cibo. E gli
unici a poterci dare un passaggio per tornare a casa», indicò verso
l'alto con il pollice, «sono i kristang. Non sappiamo nemmeno decollare
da soli. I kristang sosterranno la nostra presenza qui solo finché
saremo alleati affidabili in questa battaglia. Ci siamo offerti
volontari per venire qui, e ora che ci siamo dobbiamo fare del nostro
meglio.» Scosse la testa. «Senti, colonnello, tu e gli altri obiettori
come te ci mettete in un bel casino. Puoi chiamarci schiavi se vuoi, il
fatto è che gli umani sono subordinati ai kristang, comunque la metti.
Non siamo abituati, specie io e te, in quanto americani, a non essere
del tutto padroni dei nostri destini. Il comando Unef non sta
protestando contro la tua sorte, perché sono più preoccupati di quello
che sta succedendo a casa che di noi qui, o di un colonnello mustang.»

«Capisco.» Cercai di fare buon viso a cattivo gioco, ma in verità mi
stavo cagando sotto. Tenevo le mani unite dietro la schiena perché mi
tremavano come foglie. Non è che avessi paura di morire di fronte a un
plotone di esecuzione dei kristang, quella sorte l'avevo accettata
quando mi ero rifiutato di uccidere i criceti. Avevo paura di
comportarmi in maniera disdicevole, temevo che, quando sarebbe arrivato
il momento e sarei stato contro il muro, mi sarei pisciato addosso o
buttato a terra, piangendo e frignando per chiedere pietà, disonorando
la mia specie. Solo pensarci, pensare a quanto sarebbero state
disgustate le lucertole vedendo un essere umano comportarsi in modo così
vigliacco, mi mise un po' d'acciaio nei nervi. Avrei potuto usarlo.
Usare quel pensiero per trasformare la mia paura in odio. Odio per le
lucertole. Sì, fanculo tutta la loro specie lucertoloide puzzolente di
nazismo. Che andassero dritti all'inferno.

«Maggiore», mi spinsi via dal muro e gli offrii la mano destra, che non
stava più tremando. «Grazie di essere venuto. Nessun soldato vuole
morire da solo. Ringrazia il comando dell'Unef per me.» Ci stringemmo la
mano, avrei detto che non sapesse cosa dire. «Se mai ne avrai la
possibilità, un giorno, di' ai miei che ho fatto quella che ritenevo la
cosa giusta.»

\hfill\break

Non ebbi molto tempo per contemplare la mia sorte, perché il maggiore
Cochrane se n'era andato da meno di un minuto, quando allarmi
risuonarono in tutta la base kristang, seguiti a ruota da una tremenda
esplosione. Il pavimento della mia cella si sollevò per colpirmi in
faccia. Vedevo le stelle, mi fischiavano le orecchie. Una seconda
esplosione mi fece rimbalzare dal pavimento e una sezione della parete
esterna della cella si spaccò e crollò. La testa mi girava ancora, ma il
mio corpo non attese un invito in carta bollata, se la stava già
squagliando attraverso l'apertura prima che mi rendessi conto di cosa
stavo facendo. L'addestramento dell'esercito era stato buono, e fece
effetto proprio quando ne avevo bisogno. Per orientarmi, guardai il
cielo, un grosso errore. Proprio in quel momento, una frazione di
secondo dopo aver visto il familiare scintillio di astronavi che
saltavano in alta orbita, là in alto ci fu un'esplosione con una luce
lancinante, che mi spinse a terra, mentre macchie mi nuotavano davanti
agli occhi folgorati. Un replay di quello che era successo nella mia
città. La luce intensa doveva essere quella di una nave kristang
vaporizzata in orbita bassa. I ruhar erano tornati. Strisciai sulle
macerie del muro, mezzo cieco e sordo, e senza sapere se essere
spaventato o felice o fatalista rispetto alla svolta degli eventi. I
ruhar non erano tornati per salvare la spedizione delle Nazioni unite
dalle lucertole e nessun criceto sarebbe stato felice di vedermi. Mentre
strisciavo per terra, vista e udito ritornavano piano piano, la mia
mente correva. E ora? Dove diavolo pensavo di andare?

Che ci avessi pensato o no, che era il nocciolo della questione, il mio
addestramento militare aveva fatto effetto in pieno. Valutare la
situazione. Iniziare dai fatti di mia conoscenza. Fatto: era in
programma la mia esecuzione per mano di una specie che ora consideravo
nemica dell'umanità. Fatto: perciò, considerando il fatto numero 1,
l'autorità con cui l'Unef mi aveva arrestato e consegnato ai kristang
veniva meno, e io non avevo alcun dovere di obbedire a ordini illegali.
Fatto: acconsentendo, stavo legittimando quella merda. Fatto: questo non
rendeva meno efficace il mio ragionamento sulla legalità. Eravamo su un
pianeta alieno, intrappolati tra due specie in guerra, con poche o
nessuna speranza di tornare a casa, o anche solo sopravvivere un altro
mese. L'autorità generale dell'Unef era piuttosto limitata. Soprattutto
perché era evidente che l'accordo di resa dei ruhar e la tregua con i
kristang erano, diciamo, soggetti a interpretazione, a seconda di chi
dispiegava la flotta più grande al momento.

Le orecchie mi fischiavano ancora, ma la vista mi stava tornando, tra
una macchia e l'altra. Ero mani e ginocchia a terra fuori dalla mia
cella, per un tratto a ridosso delle macerie del muro. Attorno a me, in
tutta la base, gli edifici erano in parte o del tutto crollati. Le
esplosioni secondarie stavano ancora causando il caos. Dovevamo essere
stati colpiti da un colpo ipercinetico di cannone a rotaia, seguito a
ruota da missili intelligenti. I ruhar sapevano con esattezza dove
colpire la base kristang, dovevano aver lanciato un'ordinanza appena
usciti dal salto. Dal poco che potevo vedere, le aree carcerarie erano
quelle che avevano subito il danno minore, mentre la parte principale
del complesso kristang era più che spianata, era un cratere fumante. Un
proiettile penetratore ipersonico, arrivando anche a.05C, poteva fare
una quantità tremenda di danni. La mia ipotesi era che il dart del
cannone fosse stato seguito da vicino da un missile intelligente con
sub-munizioni, per eliminare punti critici della base che fossero
sopravvissuti all'attacco iniziale. E la mia vista era abbastanza buona
da distinguere altre luci scintillanti nel cielo, un sacco di luci
scintillanti nel cielo. Non ci furono altre esplosioni, questo mi fece
pensare che la forza simbolica di navi kristang fosse stata spazzata via
o fosse saltata via. Tutte quelle luci dimostravano che i ruhar erano
arrivati in gran numero. Mentre guardavo, ammiccando, le luci
aumentavano. Porca vacca. Non era stato solo un altro assalto. I ruhar
erano qui, in forze, per riprendersi il pianeta. La musica cambiava.

Se non fosse che avevo già preso la mia decisione, vedere segni
dell'armata ruhar l'avrebbe fatto per me. La missione dell'Unef su
Paradiso era finita, in un modo o nell'altro. Se i ruhar fossero tornati
al comando, tutti gli esseri umani sulla superficie sarebbero stati
fatti prigionieri, prima o poi. E se Paradiso era sul punto di diventare
una zona calda di battaglia, i nostri M4b1 e i mezzi aerei ruhar di cui
ci eravamo impossessati ci avrebbero soltanto intrappolati nel mezzo.

Oh, merda. Mi si accese la lampadina. Cibo. Se i kristang avessero perso
Paradiso, non ci sarebbero più state spedizioni di cibo in arrivo per
gli umani. Se avessero perso Paradiso, le fottute lucertole non
avrebbero fatto lo sforzo di rifornirci e i criceti non avrebbero avuto
accesso alla Terra per prendere cibo a uso umano. I ruhar sarebbero
stati abbastanza generosi da permettere agli esseri umani di continuare
a coltivare il loro cibo su Paradiso, sulla terra che avevano sottratto
loro? Non volevo scoprirlo.

\hfill\break

Cibo. I kristang avevano dato da mangiare una volta al giorno a me e,
presumo, agli altri prigionieri umani, quindi ci doveva essere una
scorta di cibo umano nella base e, data l'implacabile efficienza delle
lucertole, dovevano aver conservato il cibo vicino alla prigione. Dovevo
trovare quel cibo, prima di fare qualsiasi altra cosa.

La mia vista si era ripresa abbastanza perché potessi capire dove stavo
andando, l'udito era incerto, non riuscivo a dire se i suoni mi stessero
riverberando nei timpani colpiti, o provenissero da fuori dalla mia
testa. Dentro. Avevo bisogno di tornare dentro. A rigor di logica, il
magazzino del cibo doveva trovarsi dentro, non fuori dove stavo io.
Contravvenendo all'istinto, tornai con cautela nella mia ex cella e
provai ad aprire la porta. Era rotta e incastrata in un angolo, ma non
si sarebbe mossa. Di nuovo fuori, affrettandomi perché adesso avevo le
idee chiare e volevo essere lontano prima che i ruhar arrivassero o
qualsiasi sopravvissuto kristang mi cercasse, mi arrampicai lungo il
muro distrutto, cercando un modo per entrare. Fu semplice, c'era una
porta aperta. Tornato dentro, sbirciai dietro l'angolo e non vidi
nessuno in corridoio. Le porte delle celle avevano una piccola finestra
appena più in alto del mio occhio, mi misi in punta di piedi per
guardare dentro. La prima cella era vuota. La seconda anche. La porta
della terza era in parte socchiusa a causa di una grande crepa che
arrivava al soffitto e, guardando tra porta e telaio, potevo vedere i
piedi e le gambe di una persona, un umano, che giaceva sul pavimento
sotto i detriti del soffitto crollato. Aprire la porta era facile, era
chiusa solo dall'interno; tirai la maniglia e l'aprii. Ci vollero due
forti strattoni per liberare la porta dal telaio deformato e, una volta
entrato, mi fu chiaro che il mio sforzo era stato sprecato. Le gambe
erano quelle di un ufficiale francese, il suo petto era stato
schiacciato dal crollo del soffitto. L'addestramento mi aveva insegnato
che quello era il momento di cercare i sopravvissuti, non i sentimenti.

La porta della cella in fondo al corridoio era intatta, ma la parete
esterna era per la maggior parte crollata. All'interno c'era il corpo di
un maggiore dell'esercito britannico. Dall'entità del danno, dedussi che
la sua cella fosse stata colpita da un missile ruhar. Le quattro celle
successive erano vuote, poi il corridoio svoltava a destra, e dietro
l'angolo c'erano i corpi del maggiore Cochrane e del kristang che
l'aveva scortato. Erano entrambi sotto un mucchio di macerie cadute dal
soffitto. Cochrane non aveva perso molto sangue, ma non aveva polso. Il
kristang aveva il petto trafitto da quello che sembrava il pezzo di una
barra di rinforzo, notai che il sangue kristang era di un colore rosso
molto più scuro di quello umano, e mi fermai a chiedermi se avesse un
contenuto di ferro superiore. È incredibile quello che mi passa per la
testa a volte. Il maledetto kristang era disarmato, avevo sperato che la
lucertola avesse un fucile, o qualsiasi tipo di arma. Non ebbi questa
fortuna. Dovetti strisciare sopra le macerie del soffitto per superarli
e attraversare una sezione di corridoio senza porte, finché non
raggiunsi un blocco di celle. Dopo un altro paio di locali vuoti,
guardai dalla finestra e vidi una donna, una donna nera nuda. Era di
spalle, mentre cercavo di allargare una crepa nel muro esterno,
un'occhiata, per quanto rapida, mi rivelò le brutte cicatrici che aveva
sulla schiena. Con quella che poi avrei compreso essere una mossa
stupida, m'inerpicai fino al maggiore Cochrane, lo trascinai fuori dalle
macerie e gli tolsi con cura la camicia, i pantaloni e gli stivali. Mi
fece sentire irrispettoso nei suoi confronti, ma la soldatessa nuda
aveva bisogno dei vestiti di Cochrane più di quanto non ne avesse lui,
cui restavano comunque i boxer. Lo lasciai appoggiato con la schiena
contro il muro e tornai in fretta dalla donna. Sapendo che le porte
delle celle erano insonorizzate, saltai urlando e calciai la porta una
volta, poi due, poi tre, nella speranza che la prigioniera all'interno
lo prendesse per un tentativo di comunicare. Poi girai la maniglia e
spalancai la porta: «Tutto bene lì dentro? Sono Bishop, esercito
americano», dissi in un sussurro trattenuto.

«Sergente scelto Adams, signore, prima Mef¹⁸», disse con una voce
tremante di sollievo. Prima expeditionary force dei Marine. L'uniforme
del corpo dei Marine di Cochrane era appropriata per lei. Neanche avesse
una qualche importanza.

Senza guardare attorno alla porta, gettai dentro i vestiti e mi scusai
che fossero tutto quello che avevo trovato. Lei s'infilò in fretta la
camicia senza curarsi di abbottonarla, tirò la porta, l'aprì per
lanciarsi nel corridoio. «Grazie per i vestiti, ma preferirei andarmene
subito da qui.»

Merda. L'avevo vista prima come una donna e poi come un soldato, o
Marine. Era ovvio che le sarebbe interessato infinitamente di più uscire
dalla sua cella, che sapere se qualche altro soldato l'avesse vista
nuda. Guardò me, poi le mie mostrine, per un attimo confusa, perché ero
troppo giovane per essere un colonnello, poi negli occhi le balenò
l'intuizione. «Ah», riuscì a fare un mezzo saluto militare mentre
s'infilava i pantaloni, «lei è quel Bishop, il colonnello. Che diavolo
sta succedendo, signore?»

«I ruhar sono tornati e questa non è una semplice incursione, hanno
un'intera armata in cielo. Penso che si stiano riprendendo il pianeta,
nel qual caso, l'Unef resterebbe disoccupata. E\ldots{} sergente Bishop,
non colonnello. Ero un colonnello solo per via delle maledette
lucertole.»

Adams mi gettò un'occhiata che avevo visto molte volte dai sergenti
scelti: «Stronzate, signore, non può farlo».

«Fare cosa?»

«Non sono state le lucertole a nominarla colonnello, è stata l'Unef. Non
l'avrebbero fatto se non avesse portato benefici a noi umani. Il suo
grado è un vantaggio; nessun soldato rinuncia a un vantaggio sul campo
di battaglia. signore.» Di nuovo quello sguardo.

«Merda.» Cazzo, non riuscivo nemmeno a essere tutto tronfio del mio
rango. «Hai ragione, hai ragione.»

«Questi stivali sono troppo grandi, mi faranno soltanto inciampare nei
miei stessi piedi», disse Adams e li calciò via. «Andrò a piedi nudi per
ora.»

Annuii e le feci segno di seguirmi, con calma. Entrambi sobbalzammo
quando un'esplosione secondaria tuonò in tutta la base e il suono
riecheggiò in corridoio. Avvenne altre due volte, con sempre meno
violenza.

E poi da dietro l'angolo sbucò un soldato kristang.

A livello fisico, non avevo chance contro quel guerriero: ero stanco,
indebolito dallo stress e dalla fame, mentre lui era più grosso e più
forte di me e armato di fucile. A ogni altro livello, era lui a non
avere alcuna chance. Era disorientato almeno quanto me, era molto
sorpreso di trovare un essere umano vivo e fuori da una cella e si stava
guardando le spalle girando l'angolo. Avevo il vantaggio di un'enorme
scarica di adrenalina. Prima che entrambi sapessimo cosa stava
succedendo, gli strappai il fucile dalle mani ed entrai in una frenesia
omicida, alimentata dal puro terrore. Supponendo in modo inconscio che
l'arma kristang avesse una qualche caratteristica che ne impediva l'uso
non autorizzato, sbattei il calcio del fucile sotto il mento della
lucertola due volte, forse tre o più, forse colpendogli la gola, non
ricordo. Tutto quello che so è che il momento prima era in posizione
verticale e il momento dopo stava cadendo, e io gli saltavo addosso sul
pavimento, sbattendogli il calcio del fucile sul cranio più e più volte.
La mia visione fu offuscata da una nebbia rossa, in parte generata dal
sangue del kristang, in parte dal mio istinto di sopravvivenza. Impiegai
tutto, ma proprio tutto quello che avevo in corpo per ficcargli il
calcio del fucile nel cervello. Se non siete mai stati in combattimento
o coinvolti in un incidente d'auto o convinti per una frazione di
secondo di essere sul punto di morire, non potete immaginare quanto in
fretta le immagini mi passassero per la mente. Ogni briciola di odio per
la sua fottuta specie nazista, per quello che avevano fatto o stavano
facendo o progettando di fare alla Terra, ogni briciola di rabbia per
quello che aveva fatto al sergente Adams e a Miranda Collins e ad altre
donne, tutta la rabbia, trattenuta con scrupolo in nome delle regole
d'ingaggio dell'esercito americano, contro gli ignoranti fanatici
selvaggi ``religiosi'' che bruciavano le scuole e uccidevano bambini in
Nigeria, ogni stronzo che aveva cercato di fare il prepotente con me al
liceo, la mia impotenza di bambino di otto anni, mentre quella sadica
merda di Michael torturava una rana, ogni persona che mi aveva tagliato
la strada nel traffico, fino all'allenatore che aveva rimproverato il
mio giovane io impressionabile per non essere perfetto nelle partite di
baseball della Little League: tutto questo finì nel calcio del fucile,
finché non lo feci passare da parte a parte nel cranio di quel
super-guerriero geneticamente potenziato e mi ritrovai a raschiare il
cemento, o qualsiasi materiale indurito le lucertole avessero usato per
fare il pavimento.

L'altra cosa che so è che il sergente Adams mi strattonò la spalla,
cercando di riportarmi alla realtà. La squadrai con un'espressione che
deve averla terrorizzata più di quanto i kristang non avessero mai
fatto, perché si allontanò da me con un balzo all'indietro. Lasciai
cadere il fucile e mi alzai in piedi.

«Porco Cristo», disse Adams con voce roca. «Mi sto sentendo male.»

Provavo la stessa cosa. Avevo già ucciso prima, ma sempre a distanza,
sempre colpi di fucile mirati a uccidere, uccidere persone che non
sempre potevo vedere con chiarezza. Stavolta era diverso. Nella mia
rabbia, avevo schiacciato il cranio del kristang come un melone maturo,
c'erano pezzi rotti di osso bianco, il sangue rosso scuro e tutto il
resto erano come purè di patate grigio grumoso o ripieno di tacchino del
Giorno del ringraziamento, solo rivoltanti pezzi non identificabili.
«Respira», dissi, «respira a fondo e lentamente.» Non era chiaro se lo
stessi dicendo a Adams o a me stesso, perché avevo i conati di vomito.

Restammo entrambi appoggiati al muro, sopraffatti dallo shock,
respirando a fondo. Lei fu la prima a riprendere il controllo del suo
stomaco, mentre in me calava uno spaventoso picco di adrenalina. Si
avvicinò al corpo bocconi del kristang e, con mia assoluta sorpresa, gli
sferrò un calcio brutale. «Ho riconosciuto questo pezzo di merda. È lo
stronzo che mi ha torturata.» Si chinò e sputò su quella che era stata
la sua faccia. «Tutto bene, signore?»

«Ehm\ldots» Sangue kristang rosso scuro mi era schizzato addosso, anche
in faccia. Il mio mignolo sinistro era piegato e cominciò a pulsare
quando me ne accorsi, non ricordo dove si fosse slogato, o come. «Sì,
bene.»

«E adesso, signore?»

Le rivolsi uno sguardo inespressivo, pensando ancora come un sergente
semplice. Lei era un sergente scelto e mi superava per grado,
perciò\ldots{} ah, sì. Merda. Adesso ero un colonnello. Per quello che
valeva. Un colonnello in un esercito intrappolato su un pianeta in mano
al nemico. Un colonnello in un esercito che mi aveva consegnato a un
altro nemico, perché mi giustiziasse per il crimine di aver cercato di
mantenere un brandello della mia umanità. «Adams, io dico di cercare in
questo posto altri prigionieri e provviste di cibo. Poi rubiamo un
qualche mezzo di trasporto e addio.» Nel senso di ``tagliamo la corda''.

«Oorah», disse lei, che era l'equivalente Marine del corretto ``Hooah''
dell'esercito americano.

Guardai il fucile kristang che avevo tra le mani: «Meglio scoprire se
questa cosa funziona». Non sembrava aver subito alcun danno per averlo
usato come un randello contro il suo ex proprietario. Lo puntai contro
il kristang morto, me l'appoggiai alla spalla e premetti il grilletto.
Niente. «Uhm.» Un pulsante, c'era un pulsante giallo sul lato destro,
sopra il grilletto. Premendolo, il pulsante diventò rosso. Questa volta,
il grilletto funzionò, un proiettile a punta esplosiva colpì il torso
del kristang e una fontana di sangue rappreso mi schizzò addosso. Non
che me ne fossi accorto. «Giallo significa con sicura, rosso significa
senza sicura.»

«Capito.»

Trovammo soltanto altri due prigionieri vivi, la maggior parte delle
celle era vuota e, a mano a mano che ci avvicinavamo al centro della
base, aumentavano le parti di edificio crollate. Adams doveva fare
attenzione, camminando sulle macerie a piedi nudi. Il primo prigioniero
che liberammo fu un capitano donna dell'esercito indiano, un pilota di
Poiana di nome Desai. Era nuda, naturalmente, Adams entrò per
rassicurarla che eravamo soccorritori e per darle la camicia della mia
uniforme, che era abbastanza lunga da coprirle il sedere. Adams alzò le
mani e mimò gesti amichevoli, che non sarebbero stati necessari, perché
Desai parlava inglese meglio di me. Il trauma che aveva subito non aveva
influenzato la sua capacità di discernimento, non appena ebbe indossato
la mia camicia, corse fuori dalla cella dietro Adams.

L'altro prigioniero era in un certo senso un mio amico, il tenente
colonnello Chang. La porta della sua cella era stata abbattuta, ma le
pareti erano in parte crollate sopra di lui e, quando lo raggiungemmo,
era riuscito a liberarsi contorcendosi, ma aveva il piede sinistro
intrappolato sotto una sezione del soffitto. Sollevammo il pezzo rotto
abbastanza perché potesse strisciare fuori. Aveva un profondo, brutto
squarcio sul polpaccio sinistro e un taglio sulla fronte, che scosse,
dicendo che se ne sarebbe occupato più tardi. Capivo il suo punto di
vista. Oltre la cella di Chang, il corridoio era crollato. Se c'erano
umani vivi tra le macerie, non saremmo bastati noi quattro per tirarli
fuori, senza attrezzature pesanti.

«No», disse Chang, scuotendo la testa. «Le lucertole mi hanno detto che
erano rimasti solo sei prigionieri dopo la scorsa notte. Oltre a noi
quattro, c'erano un francese e un inglese.»

«Cazzo. Ho trovato un francese morto e un inglese morto laggiù. Ci siamo
tutti, allora.»

Adams mi tirò per la cintura e mi spinse dietro un angolo, tenendosi un
dito sulle labbra. «Due kristang, signore, in quell'edificio bianco, li
ho visti attraverso la porta.»

«Loro ci vedono?»

«Non penso, uno di loro era di spalle.»

Girai di scatto la testa al di sopra della mia spalla. «Torniamo
indietro, vedi se\ldots»

Tre kristang uscirono dall'edificio bianco, tutti e tre armati di
fucile, ma senza caschi o giubbotti antiproiettile. Non guardavano nella
nostra direzione, ci davano le spalle, con gli occhi rivolti al cielo.
Non esitai. O le munizioni kristang erano silenziose, oppure i loro
fucili avevano silenziatori incorporati, perché emettevano un pop-pop
solo quando i proiettili uscivano dalla canna. Un forte rumore lo fecero
le punte esplosive che colpirono i corpi dei kristang. Caddero tutti in
fretta, solo uno di loro riuscì a sparare un colpo, senza mirare a nulla
di vicino a noi. Non c'era bisogno di dare ordini, volammo verso i
kristang, prendemmo i fucili e ci nascondemmo all'interno dell'edificio
bianco. «Rosso significa senza sicura», dimostrò Adams.

«Uno di loro è ancora vivo», avvisò Chang, e alzò il fucile. Adams
spinse la canna di lato e scosse la testa con rabbia. «No, signore. Non
lei.» Guardò Desai e le due donne annuirono. Un kristang stava
dondolando da una parte all'altra a terra, avevo mancato il torso e gli
avevo fatto saltare un braccio. Desai gli ficcò due colpi in testa, poi
mormorò qualcosa in hindi, diedi per scontato che fosse del tipo
``Adios, figlio di puttana'' o qualunque fosse l'equivalente
dell'esercito indiano.

«Grazie», disse a Adams. Se Chang era arrabbiato con il sergente per
averlo redarguito, non lo diede a vedere. Sapeva che le donne erano
state torturate dalle lucertole, capiva che Desai aveva bisogno di un
sentimento di rivalsa. Far schizzare cervelli di lucertola nel raggio di
una dozzina di metri quadrati era un buon inizio.

«Qualcuno sa dove conservano il cibo da queste parti?», chiesi. «Le
lucertole mi davano un pasto pronto ogni mattina, devono avere cibo per
umani da qualche parte.»

«Non mangio da due giorni», disse Desai, e Adams annuì.

Cazzo, ora il mio singolo pasto al giorno sembrava un lusso. Tutti molto
propensi alla violenza, ispezionammo l'edificio bianco, sospettando che
ci potessero essere altri kristang. Desai trovò un armadio che conteneva
diciassette pacchetti di cibo pronto americano e la versione delle
razioni da campo dell'esercito cinese. Adams e Desai dissero che non
avevano fame, ma io insistetti che scartassero un pacchetto e lo
condividessero per iniziare, mangiando piano. Avevano bisogno di
energia. Non trovammo nient'altro di utile, né vestiti umani, né armi di
scorta e munizioni. C'erano i vestiti kristang, diedi a Desai i miei
pantaloni e indossai pantaloni kristang, lei s'infilò una camicia
kristang troppo grande e mi ridiede la mia. Pensai che, se avessimo
avuto problemi, la mia camicia, con il mio distintivo, sarebbe stata
utile. Gli stivali kristang erano troppo grandi per le donne, che
tagliarono i vestiti e si avvolsero le strisce attorno ai piedi. Mentre
Adams allacciava le sue scarpe di fortuna, un paio di Avvoltoi, le
astronavi d'assalto ruhar, rombarono sopra le nostre teste. Era tempo di
andarcene.

C'era un trio di veicoli kristang blindati da trasporto truppe
parcheggiato vicino alla recinzione, le portiere non erano chiuse, ma i
mezzi non ne volevano sapere di partire, e noi non avevamo il tempo di
cercare chiavi o chip di computer o qualunque cosa fosse necessaria.
Nelle vicinanze c'era un Crivee ruhar, il parabrezza era rotto, ma le
celle a combustibile erano cariche al 70\% e partì subito. Girammo
attorno al complesso per poter percorrere l'unica strada riparata da
alberi. Era probabile che la copertura della vegetazione non fungesse da
protezione, considerando che i ruhar avevano occhi nel cielo, ma in
qualche modo ci faceva sentire più al sicuro. L'ultima cosa che vedemmo
gettando lo sguardo al complesso distrutto delle lucertole fu un Dodo in
fase di atterraggio, scortato da un paio di Avvoltoi. Incrociai le dita
e sperai che ritenessero superfluo prendersi la briga di indagare su un
Crivee.

La copertura degli alberi s'interrompeva dopo un paio di chilometri e
stavamo viaggiando verso nord-ovest attraverso l'aperta prateria. La
nostra strada giunse a un incrocio e ordinai di fermarci. «Qualcuno sa
dove siamo?»

Nessuno, nemmeno un indizio. Verso sud, potemmo vedere una Balena e
alcuni Avvoltoi in discesa, diretti lontano da noi. Purché a debita
distanza da loro, andavano bene tutte le direzioni, perciò continuammo
verso nord-est. Poi incontrammo il capitano Clueless e il suo posto di
blocco. La strada passava sopra un ponte, sul lato opposto c'erano
diversi Crivee e soldati, soldati umani, per lo più forniti
dell'equipaggiamento dell'esercito americano, quindi eravamo felici di
incontrare amici. O almeno così pensavamo.

Era un'unità della polizia militare guidata da un sottotenente; dalla
varietà di mostrine e nazionalità dei presenti doveva avere riunito
tutte le persone che era riuscito a trovare. Puntarono i fucili sul
nostro Crivee mentre attraversavamo il ponte, poi ci ordinarono di
fermarci. Uscii, tenendo il fucile kristang sotto il braccio sinistro e
facendo il saluto con la mano destra: «A rapporto, tenente, ehm», lessi
con la coda dell'occhio il nome sul suo cartellino, «Rogers».

«Signore?» Il suo saluto di risposta era esitante. Indossavo le mostrine
di colonnello dell'esercito americano sulla giacca della divisa, tanto
macchiata di sangue kristang secco che il cartellino col mio nome era in
parte oscurato. Portavo pantaloni kristang, col loro caratteristico
motivo giallo e nero, e trasportavo un'arma kristang. Inoltre, ero
troppo giovane per essere un vero colonnello in qualunque esercito.

«Ci siamo imbattuti in alcuni kristang lungo la strada», indicai col
pollice destro dietro ai nostri Crivee, «e abbiamo acquisito
tatticamente questa attrezzatura. Non abbiamo mezzi di comunicazione»,
mostrai la mia cintura vuota, «qual è la situazione?»

«Avete attaccato i kristang? Signore, penso di doverla arrestare.» I
suoi uomini alzarono le armi. Le cose potevano mettersi male molto in
fretta.

«Tenente, a meno che non avessi due belle fette di prosciutto su quei
cazzo di occhi», quella frase stava diventando la mia preferita, «non
può essere sfuggito alla tua attenzione che questo pianeta è sotto una
nuova gestione, quella dei ruhar. La missione dell'Unef qui è finita.
Ogni umano su Paradiso è prigioniero, dei ruhar. Hai in mente di
arrestare tutto l'Unef e di consegnarci ai criceti?»

«Ho ricevuto degli ordini\ldots»

«Sono un cazzuto colonnello dell'esercito degli Stati Uniti a tutti gli
effetti e tu sei un sottotenente particolarmente idiota», dissi con
ardore, attento a tenere bene il fucile sotto il braccio sinistro.
«Anche il tenente colonnello Chang qui e il capitano Desai ti superano
per grado all'Unef e il sergente scelto Adams ti supera per
intelligenza. Abbassate le armi, ora, è un ordine.»

Desai puntò il fucile fuori dal finestrino, dritto alla testa di Rogers:
«Non tornerò in quella prigione». La canna del fucile tremava un po',
per rabbia, per paura o per un ancora basso livello di zuccheri nel
sangue.

«Neanch'io», concordò Adams uscendo dal Crivee. L'espressione sulla sua
faccia faceva più paura del suo fucile. «Colonnello Bishop, mi autorizza
a far saltare in aria la testa di questo stronzo se solo mi guarda in
modo strano?» Fissò Rogers negli occhi: «Ho trascorso gli ultimi cinque
giorni a farmi affamare e torturare dalle lucertole e ho davvero,
davvero voglia di sparare a qualcosa».

«Calmatevi, tutti.» La voce della ragione provenne dall'unità
raffazzonata di Rogers. «Tenente, forse dovremmo starli a sentire.»

Rogers annuì e i suoi abbassarono lentamente le armi. A un mio cenno, i
miei fecero lo stesso. «Così va meglio. Ora, sono il colonnello Joe
Bishop, sapete, il tizio di Barney.» La faccia di Rogers lasciò
intendere di sapere di cosa stessi parlando. Cazzo. Adams aveva ragione,
mai rinunciare a un vantaggio. «Noi quattro eravamo prigionieri dei
kristang, finché i ruhar non hanno colpito la base e ci hanno liberati,
per caso. Tutti gli altri prigionieri erano già stati fucilati o
torturati e impiccati dai nostri alleati, le lucertole. Ex alleati. Ora,
voi che avete mezzi di comunicazione», indicai il suo zPhone, «che cosa
avete sentito?»

«Niente, signore, le comunicazioni sono interrotte, anche il sistema di
localizzazione Gps è fuori uso. Tutto quello che abbiamo è un messaggio
registrato dai ruhar che ci dicono di ritirarci, cooperare con loro e
attendere ulteriori istruzioni. Niente dal comando Unef, niente e-mail o
messaggi di testo, ma la funzione di traduzione funziona.»

«Come siete messi con i rifornimenti?»

«Scarsi, signore. Acqua in abbondanza, ognuno ha un carico standard di
munizioni, ma non abbiamo nemmeno cibo sufficiente per un pasto per
tutti. La maggior parte di noi era diretta verso un franco a bordo
francese quando i ruhar sono tornati. Abbiamo visto il fumo in direzione
della base», indicò una sottile colonna di fumo a sud, «e il ponte sul
fiume in quella direzione è fuori uso. Ho messo insieme tutti quelli che
sono riuscito a trovare e li ho radunati qui. Abbiamo due francesi, tre
indiani e quattro cinesi con noi. Ho pensato che questo ponte fosse un
passaggio obbligato del traffico, signore.»

Era un pensiero rispettabile, in un certo senso. Mi avvicinai per
parlare in tranquillità. «Abbiamo circa una dozzina di confezioni di
pasti pronti. Le due donne con me sono state affamate e torturate dalle
lucertole, sarebbero state impiccate domani mattina.»

Gli occhi di Roger mostravano il suo shock: «In cosa diavolo ci siamo
cacciati, signore?».

«Ci siamo dentro fin sopra la testa, questo è certo. Carichiamo e\ldots»

«Astronavi!», gridò un soldato.

«Le loro o le nostre?», chiese il tenente Rogers. Non mi preoccupai di
sottolineare che nulla nel cielo era più ``nostro''.

«Non so ancora dirlo, signore», rispose il soldato con il binocolo.
«Aspetta, sono criceti! Tre, no, quattro Avvoltoi e un paio di Dodo,
sembra.»

Stavo per ordinare alla gente di mettersi al riparo, quando la coppia di
Avvoltoi si divise per affiancarci a est e a ovest, e i pod delle armi
si aprirono. Era chiaro che ci avevano visto ed erano interessati. Forse
ostili. «Abbassate tutti le armi, alzate le mani.»

«Signore?», chiese Rogers con sospetto. Avevo, dopotutto, ucciso dei
kristang e ora avevo intenzione di arrendermi ai ruhar.

«Tenente, forse, con molta fortuna, potremmo anche fare fuori uno di
quegli Avvoltoi, se il pilota è così stupido da avvicinarsi tanto. Gli
altri, però, ci spazzerebbero via come bersagli facili. Metti l'arma a
terra, allontanati e tieni le mani in alto.» Potevo vedere che non era
convinto. Eravamo nella galassia per combattere i ruhar in modo che la
gente sulla Terra non dovesse farlo. Arrendersi senza sparare un colpo
andava contro l'istinto di un soldato. «I ruhar ci hanno colpito qui
prima e i kristang li hanno respinti», gli ricordai, anche se stavolta
avrei preferito avere a che fare con i ruhar che con i kristang. «Se
perdi una battaglia puoi ancora vincere la guerra», dissi, «contro
qualunque delle due parti finiamo per combattere.»

Rogers ammise l'evidenza. I piloti dell'Avvoltoio non erano affatto
stupidi, i quattro volavano in alto, al sicuro fuori dalla portata degli
Javelin. Un Dodo atterrò, fece sbarcare una dozzina di soldati e subito
dopo decollò per girare verso sud. Mentre la truppa di criceti si
avvicinava, lo zPhone di Rogers si animò. Rogers toccò l'auricolare:
«Signore, ci ordinano di deporre le armi e radunarci a sud della
strada».

Ci attenemmo agli ordini. Non mi pareva che avessimo altra scelta.
Ordinai a un soldato di darmi il suo zPhone e procedetti a passo lento a
piedi, con le mani in alto, per parlare con i soldati nemici, o soldati
alieni, non sapevo chi fosse il nostro nemico in quel momento. Forse
entrambe le parti. «Sono il colonnello Joe Bishop, esercito degli Stati
Uniti», annunciai. A quanto pare questo non significava un cazzo per i
ruhar, che mi confiscarono lo zPhone e mi portarono nel gruppo con tutti
gli altri. Dopo essere rimasti seduti per venti minuti a terra sotto il
sole, un soldato ruhar ci si avvicinò e mi lanciò uno zPhone. «Sei tu il
colonnello Joe Bishop che ha avuto l'onore di parlare con», disse il
traduttore, «Bahturnah Lohgellia?»

Con chi? Ah, giusto. Avevo dimenticato il suo vero nome. «La
Bürgermeister?» Questo non si traduceva in maniera chiara per lui, che
diede un colpetto all'auricolare, così ci riprovai. «Il governatore
regionale. Mi sono incontrato con lei a Teskor. Ero il sergente Bishop
allora.»

«Sì. Sei lo stesso Joe Bishop?»

«Ta», annuii.

Questo evocò una conversazione animata tra i ruhar, non saprei dire se
fosse un bene o un male, una parte degli sguardi che mi lanciarono
tradivano una decisa ostilità. Forse alcuni di loro avevano amici a
bordo delle due Balene che avevamo abbattuto al Vettore. La guerra è
guerra, sì, ma nei loro panni sarei stato incazzato con me. Il soldato
che aveva chiesto come mi chiamassi teneva l'auricolare, come se
qualcuno gli stesse parlando. «Eri prigioniero dei kristang. Come fai a
essere libero ora e come hai fatto a procurarti le armi kristang?»,
chiese.

Indicai il cielo. «Le vostre navi hanno colpito la base kristang dove mi
tenevano prigioniero, sono scappato perché l'esplosione ha danneggiato
l'edificio e sono riuscito a uscire attraverso un muro spaccato.»
Toccando la mia uniforme insanguinata, aggiunsi: «Ho ucciso un kristang
spappolandogli il cranio in una poltiglia insanguinata con il suo
fucile», mimai l'azione, nel caso le mie parole non si traducessero
bene, «poi ho sparato ad altre tre lucertole con lo stesso fucile. Altri
tre umani sono fuggiti con me. Abbiamo ucciso solo quattro lucertole,
perché è tutto quello che siamo riusciti a trovare», questa mi uscì con
molta più spavalderia di quanta fosse nelle mie intenzioni. «Ho risposto
alla tua domanda?»

Era una gara a chi strabuzzava di più gli occhi, tra il soldato ruhar e
il tenente Rogers, che rimase in silenzio. «Hai ucciso dei kristang?»,
chiese il criceto.

«I kristang giustiziavano gli umani e stavano per giustiziare noi»,
indicai i miei tre ex compagni di prigionia, «perché ci siamo rifiutati
di obbedire all'ordine di uccidere dei civili ruhar. Le lucertole hanno
picchiato e torturato le nostre femmine, le donne», aggiunsi, per
sottolineare che le donne erano persone, a dispetto di quello che ne
pensavano le lucertole.

Il soldato ruhar ebbe una conversazione animata con chiunque stesse
parlando nel suo auricolare, poi con altri due soldati. Era una sorta di
litigio, è probabile che il ruhar a terra stesse bisticciando con
qualche stupido fobbit ruhar, che se ne stava seduto sano e salvo in
orbita.

C'era da capirlo. Infine, il soldato ruhar che aveva parlato con me mi
fece un cenno: «Tu e gli altri ex prigionieri dei kristang verrete con
me». Doveva aver chiamato il Dodo, perché si avvicinava e scendeva in
fretta.

«Alt, alt, aspetta», protesi la mano, con il palmo in avanti, in quello
che speravo fosse un altro gesto del linguaggio universale del corpo.
«Non vado da nessuna parte finché non so cosa ne sarà della mia gente
qui.» Rogers e io ci scambiammo uno sguardo. Essendo l'ufficiale di più
alto grado, erano sotto la mia responsabilità, anche se a stento li
conoscevo. In quanto unità raffazzonata, non è che fossero poi da molto
tempo agli ordini di Roger.

Un breve lampo d'irritazione, che di sicuro rientrava nel linguaggio
universale del corpo, percorse la faccia del criceto, poi annuì. Da
soldato a soldato, poteva comprendere la mia preoccupazione: «Saranno
scortati fino a un punto di raccolta a ovest, dove stiamo distribuendo
cibo umano. Quello che accadrà dopo dipende dall'esito delle discussioni
tra i miei capi e i tuoi».

``E dal fatto che i kristang tornino o meno'', pensai, ma non lo dissi a
voce alta.

Viaggiare in un Dodo ruhar non era molto diverso dal viaggiare in una
navicella kristang, non fosse che questa volta ero prigioniero di
guerra. Ora che ci penso, forse prima eravamo tutti prigionieri di
guerra dei kristang, solo che ancora non lo sapevamo. I ruhar avevano
preso le nostre armi, ma avevo insistito che Desai e Adams portassero
con sé due pasti pronti per ciascuna, in modo da recuperare le forze.
Avevano entrambe tentato di protestare, perché non volevano un
trattamento speciale, fino a quando non avevo detto loro che quello di
mangiare era un ordine. Io e Chang eravamo affamati, ma avevamo fatto
almeno una colazione striminzita quella mattina. Pochi minuti dopo che
il Dodo era tornato rombando nel cielo, un membro dell'equipaggio ruhar
si avvicinò al mio posto e mi consegnò uno zPhone ruhar. Misi
l'auricolare. «Buongiorno.»

«Colonnello Bishop, buongiorno a lei. Sa con chi sta parlando?»

Riconobbi subito la voce stridula. Cazzo, perché non riuscivo a
ricordare il suo nome? Lahtoodah qualcosa? «Con la Bürgermeister,
signora.»

«Sì.» Sembrava divertita. «Mi fa piacere sentire che è sopravvissuto
alla prigionia dei kristang. Ho protestato con i suoi capi quando lei e
gli altri prigionieri di coscienza siete stati arrestati. Sapevamo per
esperienza quello che i kristang avrebbero fatto.»

Non fece parola del piccolissimo, cavilloso dettaglio: il motivo per cui
ero diventato prigioniero di coscienza era che i ruhar avevano attaccato
le truppe Unef in tutto il pianeta. Mentre erano ancora in vigore una
tregua e un accordo per evacuare il pianeta. Così, attacchi a
tradimento. Comunque anch'io non ne feci parola. «Grazie.» Mia madre
sarebbe stata orgogliosa di quanto fossi educato. «Che ne sarà degli
umani su», mi sforzai di ricordare il nome ruhar del pianeta, «Gehtanu?»
Segnate un punto a mio favore.

«Stiamo negoziando con i vostri capi. Ulteriori terreni saranno
accantonati per far crescere colture di alimenti umani. La sua gente
verrà sistemata in numerosi ampi campi\ldots»

«Pessima idea. Non fatelo.» Non l'avevo mai interrotta prima.

«Fare cosa?»

«Concentrare gli umani in campi. Abbiamo ceduto questo pianeta, i
kristang ci considereranno traditori. Se torneranno indietro, anche solo
per un'incursione, un gran numero di umani in un solo posto sarà un
grande bersaglio succulento per i cannoni a rotaia.»

«Non ci avevo pensato.»

Sono certo che il comando dell'Unef ci stava pensando. «Ci aiuterete a
coltivare il nostro cibo, ma non ci saranno più rifornimenti dalla
Terra?»

«È improbabile che i kristang faranno qualsiasi sforzo per rifornirvi, a
meno che non si aspettino di riprendere questo pianeta. E questo non
accadrà. Hanno subito una devastante sconfitta a opera delle nostre
forze nello spazio, ecco perché siamo in grado di prendere di nuovo il
controllo qui. Posso dirvi ora che la nostra invasione fallita, quando
avete difeso il Vettore, era una finta, intesa ad attirare un nutrito
gruppo di battaglia thuranin e kristang in quell'area. Ha funzionato, e
le nostre forze hanno distrutto quel gruppo di battaglia oggi. La nostra
intelligence riferisce che i thuranin non pensano più che valga la pena
combattere per questo pianeta e non sosterranno gli sforzi fatti dai
kristang per tornare. È probabile che le forze nemiche rimaste nella
zona attueranno attacchi e incursioni per darci filo da torcere qui, ma
non sono in grado di organizzare una campagna di flotta duratura in
questo momento.»

Merda. Quindi, i miei sforzi per difendere il Vettore erano stati tutti
inutili. E tutti i soldati ruhar su quelle due Balene e gli umani nella
sala mensa erano morti per niente. «In tal caso, vi suggerisco di
spostare tutti gli esseri umani nel piccolo continente a sud, si chiama
Lemuria», lo pronunciai lentamente perché sapevo che era intraducibile,
«sulle mappe umane. Possiamo stabilirvi piccoli insediamenti, incentrati
sulle fattorie. Non ci sono molte persone lì, penso che sia meglio
tenere umani e ruhar separati.»

«Anche di questo si sta discutendo. Speriamo che la vostra gente qui
possa unirsi a noi, alla fine.»

«Cambiare parte nella guerra? Accetteremo una tregua, ma non passeremo
da una parte all'altra, non finché i kristang controlleranno la Terra.»

«Colonnello Bishop, di sicuro vede\ldots»

«Vedo che la nostra Expeditionary force umana è bloccata qui, in un
posto dove non possiamo mangiare il cibo locale, e vedo che non ci
saranno più rifornimenti dalla Terra. Vedo che voi e i kristang
continuerete a combattere per questo pianeta e noi siamo presi in mezzo.
Vedo che le maledette lucertole hanno il controllo del mio pianeta
natale, e vedo che niente di quello che faccio qui farà alcuna
differenza a casa. È quello che vedo. Se i ruhar sarebbero alleati più
onorevoli dei kristang o no non è rilevante al momento.»

Ci fu una pausa lunga prima che rispondesse, in parte perché il
traduttore dovette mettersi in pari col mio discorso piuttosto esteso.
«Capisco che si trova in una posizione difficile. Colonnello Bishop,
quando le trattative saranno concluse e la situazione sarà più stabile,
vorrei che prendesse in considerazione l'eventualità di venire a
lavorare con me, come mio collegamento.»

Era tradotto bene? ``Collegamento''?

«Come forse avrà intuito, non sono soltanto un governatore regionale.
Sono quello che potremmo chiamare il viceamministratore di questo
pianeta.»

Porca puttana. No, non ne avevo idea. Ebbi la prontezza di non
ammetterlo. ``Perché io? Ci sono un sacco di umani che hanno svolto
incarichi di collegamento.» O, meglio, davo per scontato che l'Unef
avesse persone di collegamento. «Ho ucciso i soldati ruhar al Vettore,
molti tra la sua gente devono odiarmi.»

«È un peccato, sì. Ha anche trattato bene un soldato ruhar catturato sul
suo pianeta, però, e di recente ha rischiato la sua vita qui per
proteggere i civili ruhar. I soldati li ha uccisi in combattimento.
L'avrebbero uccisa. Non conosco molti umani, conosco lei e credo sia una
persona di carattere.»

«Grazie. Prenderò in considerazione la sua offerta. Se i miei capi sono
d'accordo, mi capisce.»